
\section{Learning the regularizer}

Assume, the regularizer depends on some parameter $\theta$. Then learning it is equivalent to minimizing the negative log-likelihood:
\begin{align*}
	\min_\theta \mathcal{L}(\theta) = \lambda \int_{f} R_\theta(f)d\pi(f) + \log Z(\theta)
\end{align*}
where  $Z(\theta) = \int_{f \in X} \exp(-R_\theta(f)) \nu(df)$ is a normalizing constant, $X$ is function space, $\nu$ base reference measure on $X$.

Or, if we are given the empirical observations of functions $\{ f_i \}_{i=1}^t$ then
\begin{align*}
	\min_\theta \mathcal{L}(\theta) =  \frac{\lambda}{t} \sum_{i} R_\theta(f_i) + \log Z(\theta)
\end{align*}

Assuming $R_\theta(f) = \langle \theta, \Phi(f) \rangle$ (linear in $\theta$) out problem is convex, yaaaaay! Parameter $\theta$ can be here an infinite-dimensional measure or a vector, hence, the minimization problem can be solved by gradient descent on $\theta$.

The gradient of the data term is
\[
	\nabla_\theta
	\left(
	\frac{1}{T}\sum_{t=1}^T \langle \theta,\Phi(f^{(t)})\rangle
	\right)
	=
	\frac{1}{T}\sum_{t=1}^T \Phi(f^{(t)})\eqqcolon \mathbb E_{\mathrm{data}}\!\left[\Phi(f)\right].
\]

For the partition function term, differentiation under the integral yields
\begin{align*}
	\nabla_\theta \log Z(\theta)
	 & =
	\frac{1}{Z(\theta)} \nabla_\theta Z(\theta)             \\
	 & =
	\frac{1}{Z(\theta)}
	\int \left(-\Phi(f)\right)
	\exp\!\big(-\langle \theta,\Phi(f)\rangle\big)\,\nu(df) \\
	 & =
	- \mathbb E_{f \sim p_\theta}\!\left[\Phi(f)\right].
\end{align*}

Combining the two terms, the gradient of the objective is
\[
	\boxed{
		\nabla_\theta \mathcal L(\theta)
		=
		\mathbb E_{\mathrm{data}}\!\left[\Phi(f)\right]
		-
		\mathbb E_{f \sim p_\theta}\!\left[\Phi(f)\right],
	}.
\]



A gradient descent step with step size $\eta>0$ is therefore given by
\[
	\boxed{
		\theta_{k+1}
		=
		\theta_k
		-
		\eta
		\left(
		\mathbb E_{\mathrm{data}}\!\left[\Phi(f)\right]
		-
		\mathbb E_{f \sim p_{\theta_k}}\!\left[\Phi(f)\right]
		\right).
	}
\]

If $\theta$ is constrained (e.g., $\theta \ge 0$ or $\theta$ lies in a simplex), the update is followed by projection onto the feasible set $\Theta$:
\[
	\theta_{k+1}
	=
	\Pi_\Theta
	\left[
		\theta_k
		-
		\eta
		\left(
		\mathbb E_{\mathrm{data}}\!\left[\Phi(f)\right]
		-
		\mathbb E_{f \sim p_{\theta_k}}\!\left[\Phi(f)\right]
		\right)
		\right].
\]

At a stationary point $\theta^\star$, the optimality condition is the moment-matching identity
\[
	\mathbb E_{\mathrm{data}}\!\left[\Phi(f)\right]
	=
	\mathbb E_{f \sim p_{\theta^\star}}\!\left[\Phi(f)\right].
\]

$\Phi(f)$ here a statistic extractor on function space.

\section{Examples}

Note, in all the examples that follow $\Phi(f)$ is convex in $f$.
\paragraph{Example 1 (Energy regularization).}
\[
	\Phi : L^2(M) \to \mathbb{R},
	\qquad
	\Phi(f) := \|f\|_{L^2(M)}^2.
\]

The induced regularizer $R(f)=\|f\|_{L^2}^2$ corresponds to ridge/Tikhonov regularization.
This choice measures only the magnitude of $f$ and ignores geometric structure.


\subsection{First-Order Smoothness}

\paragraph{Example 2 (Sobolev $H^1$ regularization).}
\[
	\Phi : H^1(M) \to \mathbb{R},
	\qquad
	\Phi(f) := \int_M \|\nabla f(x)\|^2 \, dx.
\]

This feature measures global smoothness.
The null space consists of constant functions.


\paragraph{Example 3 (Anisotropic Sobolev regularization).}
Let $\{Q_j\}_{j=1}^k$ be fixed symmetric positive semidefinite matrices.
Define
\[
	\Phi : H^1(M) \to \mathbb{R}^k,
	\qquad
	\Phi_j(f) := \int_M \langle \nabla f(x), Q_j \nabla f(x)\rangle \, dx.
\]

Then
\[
	R_\theta(f) = \sum_{j=1}^k \theta_j \Phi_j(f)
\]
corresponds to learning a Riemannian metric.


\subsection{Total Variation and Sparsity}

\paragraph{Example 4 (First-order total variation).}
\[
	\Phi : BV(M) \to \mathbb{R}_+,
	\qquad
	\Phi(f) := \|\nabla f\|_{\mathcal M}.
\]

This feature measures the total mass of the gradient as a Radon measure and promotes piecewise-constant functions.


\subsection{Second-Order and Curvature-Based Statistics}\label{par:anisotropic-second-order-tv}

\paragraph{Example 5 (Second-order total variation / splines).}
\[
	\Phi : X \to \mathbb{R}_+,
	\qquad
	\Phi(f) := \|\Delta f\|_{\mathcal M},
\]
where $X$ is the space of functions with measure-valued Laplacian.

The null space consists of affine functions, and minimizers are adaptive splines.


\paragraph{Example 6 (Anisotropic second-order TV).}
Let $\mathcal S_+$ denote the cone of symmetric positive semidefinite matrices.
Define
\[
	\Phi : X \to \mathcal C(\mathcal S_+,\mathbb{R}_+),
	\qquad
	\Phi(f)(Q) := \|\Delta_Q f\|_{\mathcal M},
\]
where $\Delta_Q f := \nabla\cdot(Q\nabla f)$.

Here:
\begin{itemize}
	\item domain: functions with measure-valued anisotropic Laplacian,
	\item codomain: nonnegative functions on $\mathcal S_+$,
	\item regularizer:
	      \[
		      R_\theta(f) = \int_{\mathcal S_+} \|\Delta_Q f\|_{\mathcal M}\, d\theta(Q).
	      \]
\end{itemize}

This promotes sparse, geometry-aligned curvature and induces a Finsler total variation.

\subsection{Spectral Statistics}

\paragraph{Example 7 (Spectral energy features).}
Let $\{\psi_k\}$ be eigenfunctions of a Laplacian on $M$.
Define
\[
	\Phi : L^2(M) \to \ell^1,
	\qquad
	\Phi_k(f) := |\langle f,\psi_k\rangle|^2.
\]

The corresponding regularizer weights energy by frequency and induces a spectral bias.


\subsection{Measure-Based and Neural Representations}

\paragraph{Example 8 (Barron / ReLU representation).}
Let
\[
	f_\mu(x) = \int_{\mathbb S^{d-1}} \sigma(\langle w,x\rangle)\, d\mu(w),
\]
with $\mu$ a finite signed Radon measure.
Define
\[
	\Phi : \mathcal B \to \mathcal M_+(\mathbb S^{d-1}),
	\qquad
	\Phi(f_\mu) := |\mu|.
\]

Here:
\begin{itemize}
	\item domain: Barron space $\mathcal B$,
	\item codomain: finite nonnegative measures,
	\item regularizer: $R(f_\mu)=\|\mu\|_{\mathrm{TV}}$.
\end{itemize}

This promotes sparse ridge representations.


\paragraph{Example 9 (Finsler--TV / Radon--Laplacian).}
Combining the previous examples,
\[
	\Phi : \mathcal B \to \mathcal C(\mathcal S_+,\mathbb{R}_+),
	\qquad
	\Phi(f_\mu)(Q) := \int \sqrt{w^\top Q w}\,|\mu|(dw).
\]

Equivalently,
\[
	R(f_\mu)
	=
	\int_{\mathcal S_+} \Phi(f_\mu)(Q)\, d\theta(Q)
	=
	\int F(w)\,|\mu|(dw),
	\quad
	F(w):=\int \sqrt{w^\top Q w}\, d\theta(Q).
\]

This is a Finsler total variation norm on kink directions.



\section{Back to the problem}
Given all distribution over functions minimizing the negative log-likelihood is
\begin{align*}
	\min_\theta \mathcal{L}(\theta) =
	% \underbrace{
	\lambda \int_{f} R_\theta(f)d\pi(f)
	% }_{\text{energy of functions}}
	+
	\log Z(\theta)
\end{align*}


where  $Z(\theta) = \int_{f \in X} \exp(-R_\theta(f)) \nu(df)$ is a normalizing constant, $X$ is function space, $\nu$ base reference measure on $X$.

However, everything we observe is just a discrete sample. Hence, given the empirical observations of functions $\{ f_i \}_{i=1}^t$ our minimization looks like
\begin{align*}
	\min_\theta \mathcal{L}(\theta) =  \frac{\lambda}{t} \sum_{i} R_\theta(f_i) + \log Z(\theta)
\end{align*}

Again, \textbf{everything} we observe is a discrete sample. Usually, we have access to each function $f$ via its values: $D = \{(x_i, y_i)\}$ where $y_i = f(x_i) + \mathcal{N}(0, \sigma^2)$. Note, that
\begin{align*}
	\log p_{\theta}(D_Y | f, D_X) \propto - \frac{1}{\sigma^2}\sum_{(x, y) \in D} \|y - f(x)\|_2^2
	= \ell(D, f)
\end{align*}
Hence,

\begin{align*}
	-\log p_\theta\left(D_Y \mid D_X\right)=-\log \int \exp \left(-\ell(D, f)-\lambda R_\theta(f)\right) \nu(d f)+\log Z(\theta)
\end{align*}

So given datasets $\{ D^{(i)} \}_{i=1}^t$ our optimization problem is
\begin{align}\label{eq:loss-w-data}
	\min_\theta \mathcal{L}(\theta) =
	-\frac{1}{t} \sum_{i=1}^t \log \int \exp \left(-\ell\left(D^{(i)}, f\right)-\lambda R_\theta(f)\right) \nu(d f)+\log Z(\theta)
\end{align}

However, the term $\log \int \exp \ldots $ is intractable :( This is why the masking comes into play in practice.

\section{Masking}

% Assume:
% \begin{itemize}
%     \item The dataset $D$ is large (large $n$ )
%     \item  Noise variance $\sigma^2$ is small
%     \item The regularizer $R_\theta$ is coercive
% \end{itemize}

% Then the functional
% $$
% E_\theta(D, f):=\ell(D, f)+\lambda R_\theta(f)
% $$

% has a unique minimizer
% $
% \hat{f}_\theta:=\arg \min _f E_\theta(D, f)
% $.
% Around $\hat{f}_\theta$, the energy grows quadratically. Hence, the posterior over $f$,
% $
% p_\theta(f \mid D) \propto e^{-E_\theta(D, f)},
% $
% is sharply concentrated near $\hat{f}_\theta$


\subsection{Posterior concentration around the MAP estimator}
\label{sec:posterior-concentration}

Fix a parameter $\theta$ and a dataset $D=\{(x_j,y_j)\}_{j=1}^n$.
Define the energy functional
\[
	E_\theta(f)
	:=
	\ell(D,f)+\lambda R_\theta(f),
\]
where $\ell(D,f)$ denotes the empirical negative log-likelihood and
$R_\theta$ is a regularizer.

Assume the following:
\begin{enumerate}
	\item The dataset $D$ is large (large $n$) and noise variance $\sigma^2$ is small
	\item (\emph{Coercivity}) The functional $R_\theta$ is coercive on $X$
	      modulo a finite-dimensional nullspace $\mathcal N$, i.e.,
	      $\|f\|_X\to\infty$ with $f\perp\mathcal N$ implies $R_\theta(f)\to\infty$.
	\item (\emph{Strict convexity}) The functional $\ell(D,\cdot)$ is strictly
	      convex on $X/\mathcal N$.
	\item (\emph{Regularity}) The map $f\mapsto E_\theta(f)$ is twice
	      differentiable in a neighborhood of its minimizer.
\end{enumerate}

Under these assumptions, the minimization problem
\[
	\hat f_\theta
	\in
	\arg\min_{f\in X} E_\theta(f)
\]
admits a unique solution modulo $\mathcal N$.
Moreover, the posterior distribution
\[
	p_\theta(f\mid D)
	\propto
	\exp\!\big(-E_\theta(f)\big)
\]
concentrates in probability around $\hat f_\theta$ as either
$n\to\infty$ or $\sigma^2\to 0$.
In particular, for any neighborhood $U$ of $\hat f_\theta$ in $X$,
\[
	p_\theta(U\mid D)\;\longrightarrow\;1.
\]


\subsection{Laplace approximation of the marginal likelihood}
\label{sec:laplace}

Let $E_\theta(f)=\ell(D,f)+\lambda R_\theta(f)$ and let
$\hat f_\theta$ denote its unique minimizer.
Define the Hessian operator
\[
	H_\theta := \nabla^2 E_\theta(\hat f_\theta),
\]
acting on the tangent space of $X/\mathcal N$.

By a second-order Taylor expansion,
\[
	E_\theta(f)
	=
	E_\theta(\hat f_\theta)
	+
	\frac{1}{2}
	\langle f-\hat f_\theta,\,
	H_\theta (f-\hat f_\theta)\rangle
	+
	o(\|f-\hat f_\theta\|_X^2).
\]

Substituting this expansion into the marginal likelihood integral yields
the Laplace approximation
\begin{equation}
	\label{eq:laplace-approx}
	\int_X
	\exp\!\big(-E_\theta(f)\big)\,\nu(df)
	\;\approx\;
	\exp\!\big(-E_\theta(\hat f_\theta)\big)
	\int
	\exp\!\left(
	-\tfrac{1}{2}
	\langle h, H_\theta h\rangle
	\right)\,dh.
\end{equation}

The remaining integral is Gaussian and evaluates (up to a multiplicative
constant) to $(\det H_\theta)^{-1/2}$, where $\det H_\theta$ denotes a suitable
regularized determinant. Consequently,
\begin{equation}
	\label{eq:laplace-result}
	\int_X
	\exp\!\big(-\ell(D,f)-\lambda R_\theta(f)\big)\,\nu(df)
	\;\approx\;
	\exp\!\big(-\ell(D,\hat f_\theta)-\lambda R_\theta(\hat f_\theta)\big)
	(\det H_\theta)^{-1/2}.
\end{equation}


\subsection{Decomposition of the marginal likelihood objective}
\label{sec:objective-decomposition}

Substituting the Laplace approximation~\eqref{eq:laplace-result} into the
negative log marginal likelihood yields
\begin{equation}
	\label{eq:laplace-objective}
	-\log p_\theta(D_Y\mid D_X)
	\;\approx\;
	\ell(D,\hat f_\theta)
	+
	\lambda R_\theta(\hat f_\theta)
	+
	\frac{1}{2}\log\det H_\theta
	-
	\log Z(\theta),
\end{equation}
up to additive constants independent of $\theta$.

This expression decomposes the marginal likelihood into four interpretable
components:
\begin{enumerate}
	\item \emph{Data-fit term} $\ell(D,\hat f_\theta)$, measuring how well the
	      MAP estimator explains the observations;
	\item \emph{Regularization cost} $\lambda R_\theta(\hat f_\theta)$,
	      encoding alignment between the learned function and the inductive bias
	      imposed by $\theta$;
	\item \emph{Local complexity (Occam) penalty}
	      $\frac{1}{2}\log\det H_\theta$, which penalizes flat or unstable directions
	      around the MAP solution;
	\item \emph{Global complexity penalty} $-\log Z(\theta)$, measuring the
	      overall volume of functions favored by the prior.
\end{enumerate}

Together, these terms characterize the tradeoff between data fit,
regularization, and model complexity in the marginal likelihood objective.


\section{Back to the optimization}
Given the collection of datasets $\mathcal{D} = \left\{\left(D^{(i)}_{\text{train}}, D^{(i)}_{\text{test}}\right)\right\}_{i=1}^t$
Substituting this into \Cref{eq:loss-w-data} we have

\begin{gather*}
	\min_\theta \mathcal{L}(\theta) =
	\frac{1}{t} \sum_{i=1}^t
	\ell\left(D_{\text{test}}^{(i)},\hat f^{(i)}_\theta\right)
	+
	\lambda R_\theta\left(\hat f^{(i)}_\theta\right)
	+
	\frac{1}{2}\log\det H_\theta
	-
	\log Z(\theta)
	\\
	\text{where } \hat f_\theta^{(i)} = \argmin_f
	\ell\left(D_{\text{train}}^{(i)},  f\right)
	+
	\lambda R_\theta\left(f\right)
\end{gather*}

This is bilever optimization. Even though it lost the convexity, it is tractable and can be solved in practice by alternating gradient descent. I claim that gradient descent of the outer problem  = training a transformer, gradient descent of the inner problem is a transformer forward pass.

\section{Inner problem}

For a dataset $D = \{(x_i, y_i)\}$ we have

\begin{align*}
	\hat f_\theta = \argmin_f
	\ell\left(D,  f\right)
	+
	\lambda R_\theta\left(f\right)
	=
	\argmin_f
	\| f(x_i) - y_i\|_2^2
	+
	\lambda \langle \theta, \Phi(f) \rangle
\end{align*}

Note, that if $\Phi$ is convex in $f$ then this can be solved by subgradient descent on $f$:
\begin{align*}
	f^{(t+1)} = f^{(t)} - \eta \left[ 2\sum_i \bigl(f^{(t)}(x_i) - y_i\bigr) \delta_{x_i} + \lambda \, g^{(t)} \right]
\end{align*}
where $g^{(t)} \in \partial_f R_\theta(f^{(t)})$ is a subgradient of the regularizer.

\subsection{How transformer forward pass can implement a (sub)gradient descent}

Now lets go through some of the examples in \Cref{par:anisotropic-second-order-tv}
to see how transformer-like networks can implement a (sub)gradient $g$ descent for the inner problem.

\paragraph{Energy regularization}

\[
	\Phi : L^2(M) \to \mathbb{R},
	\qquad
	\Phi(f) := \|f\|_{L^2(M)}^2 \qquad R_\theta(f) = \theta \|f\|_{L^2(M)}^2
	\qquad
	g = 2f
\]

\paragraph{Anisotropic Sobolev Regularization}
Parameter $\theta$ is a measure over PSD matrices.
\[
	\Phi(f)(Q) := \| \nabla f(x) \|^2_{Q} \qquad R_\theta(f) =  \int_{Q \in \mathbf{S}_+}\| \nabla f(x) \|^2_{Q}  d \theta(Q) \qquad g = - \int_{Q \in \mathbf{S}_+} \Delta_Q  f  d \theta(Q)
\]


\paragraph{Anisotropic second-order TV}

\[
	\Phi : X \to \mathcal C(\mathcal S_+,\mathbb{R}_+),
	\qquad
	\Phi(f)(Q) := \|\Delta_Q f\|_{\mathcal M},
	\qquad
	g = \int_{\mathcal S_+}  \Delta_Q \sign(\Delta_Q f) d\theta(Q)
\]
where $\Delta_Q f := \nabla\cdot(Q\nabla f)$.


As the functions are defined on a "flat" domain $\domain \subset R^{d}$, the discrete version of $\Delta_Q$ looks like:
\begin{align*}
	L \in R^{n \times n}: L_{ij} = \exp\left(\frac{-\|x_i - x_j\|_{Q^{-1}}^2}{\epsilon} \right)
\end{align*}

If our domain is curved manifold $\domain$ we can W.L.O.G. (by Takahashi's thm hehe) assume that it's $\domain \in S^{d-1}$, hence, the discrete version of $\Delta_Q$ will require projection $P_i = I - x_i x_i^\top$:
\begin{align*}
	L \in R^{n \times n}: L_{ij}
	 & = \exp\left(\frac{-(x_i - x_j)^\top P_i Q^{-1} P_i  (x_i - x_j)^\top}{\epsilon} \right)                                                               \\
	 & = \exp \left(-\frac{\left(x_j-\left(x_i^{\top} x_j\right) x_i\right)^{\top} Q^{-1} \left(x_j-\left(x_i^{\top} x_j\right) x_i\right)}{\epsilon}\right)
\end{align*}

Note, that all the exp terms are kind of similar to the original transformer softmax terms. I think, softmax should be replaced by those. The final idea is that a transformish NN with infinite amount of heads approximates the distribution over $Q$.